\doublespacing % Do not change - required

\chapter{Professional Issues and Project Management}
\label{chLast}

%%%%%%%%%%%%%%%%%%%%%%%%%%%%%%%%%%%%%%%
% IMPORTANT
\begin{spacing}{1}
\minitoc
\end{spacing}
\thesisspacing
%%%%%%%%%%%%%%%%%%%%%%%%%%%%%%%%%%%%%%%

\section{Project Plan}

The project was organised into work packages (WPs) covering literature review, baseline reproduction, attack integration, HRAC implementation, experimental evaluation, and final write-up. Early effort was concentrated on establishing a reliable experimental pipeline, including datasets, training loops, attack baselines, logging, and plotting, so that later conclusions would not be limited by tooling or reproducibility issues.

Progress was managed through weekly milestones defined by measurable outcomes rather than only dates. Core milestones included reproducing a stable benign baseline, validating each attack setting, integrating HRAC into the common training loop, and producing result tables, ablations, and diagnostic logs. When outcomes deviated from expectations, the plan was revised by updating estimates and re-prioritising tasks, particularly during phases where debugging and experimental variance dominated progress.

Risk management focused on training instability under non-IID partitions, hyperparameter sensitivity, and limited feedback cycles. Mitigations included staged evaluation from IID to non-IID settings, fixed configurations for main comparisons, and early production of intermediate results. A contingency path was also defined: if HRAC did not consistently outperform strong baselines, the project would still deliver a defensible contribution through ablation studies, failure-mode analysis, and clear reporting of limitations.

\section{Evaluation Plan}

The evaluation was structured around three questions: whether HRAC preserves benign utility, whether it improves robustness under Byzantine behaviour without oracle attacker labels, and whether its logged residual and weighting signals help explain the observed behaviour.

Results are reported across IID and non-IID partitions, with the main CIFAR-10 attack settings using $b=3$ malicious clients out of $N=10$. The attack set includes protocol-compliant label flipping, structural model manipulation, history-aware manipulation, and variance- or direction-based attacks. All methods are evaluated under matched training settings, and the limited seed and attacker-fraction coverage is reported explicitly in Chapter~\ref{ch:experiments-results}.

Success is measured primarily through test-accuracy trajectories, final and best accuracy, and component ablations. Diagnostic logs, including residual thresholds, $\mu/\nu$ histories, normalised weights, and attacker-aware summaries, are used to support interpretation rather than to define success alone.

Component-level justification is provided through ablation studies that disable individual HRAC mechanisms while keeping the surrounding pipeline fixed. Reproducibility is supported by cached records, logged configurations, and standardised plots and tables across methods. Additional seeds, datasets, and attacker fractions remain important extensions beyond the current scope.

\section{Legal and Ethical Matters}

This project investigates robustness in federated/distributed learning under the Byzantine threat model. The work is conducted as a research and engineering study using established open-source software and publicly available benchmark datasets for academic evaluation. It does not involve deployment in a live operational environment or the collection of new personal data.

\textbf{Legal considerations.}
The implementation relies on third-party libraries and datasets distributed under open licences. Their licence terms and attribution requirements are respected through appropriate citation and acknowledgement in the dissertation and accompanying codebase. From a data protection perspective, experiments are performed on standard public datasets rather than personal or sensitive user records. The project nevertheless recognises that, in real-world federated learning deployments, model updates may carry information that could be considered sensitive. Any future deployment would therefore require suitable governance and technical controls (e.g., access control, secure aggregation, and procedures consistent with applicable data protection regulation such as GDPR).

\smallskip

\textbf{Ethical considerations.}
The project includes adversarial behaviours to evaluate robustness systematically; the intent is defensive and scientific rather than enabling misuse. Ethical risks are mitigated by focusing on aggregate results and interpretable diagnostic statistics, and by avoiding unnecessary dissemination of operational misuse guidance. The work also considers the ethical risk of unfairly penalising benign clients under non-IID data heterogeneity. To reduce such harm, HRAC compares clients mainly with their own histories and applies continuous suppression rather than rigid exclusion.

\section{Impact}

The main environmental impact of this project comes from compute used for repeated training runs (attacks, baselines, and ablations). This is limited by keeping experiments tightly scoped, reusing standard benchmarks, and logging results carefully to avoid unnecessary reruns. The defence itself is lightweight on the server side because it uses vector norm operations and per-client EMA states, so it does not add substantial overhead per round.

Societally, more robust federated learning can make privacy-preserving systems more reliable in areas such as edge sensing, healthcare, and industrial monitoring. The main risks are increased compute in large deployments and the possibility that benign but highly non-IID clients are down-weighted over time. These are reduced by using client-specific histories, soft weighting rather than hard exclusion, and participation diagnostics to catch unwanted behaviour early.

\section{EDI}

Equality, Diversity, and Inclusion considerations are reflected in the project's emphasis on heterogeneous client behaviour. Non-IID data is treated as a first-class condition, which in real deployments may correspond to diversity in users, devices, environments, and usage patterns. Robust aggregation methods can unintentionally disadvantage minority or atypical yet benign clients if they rely on tight clustering assumptions. HRAC mitigates this risk by using each client's historical residual behaviour as the main reference, aiming to remain tolerant to legitimate variability while still resisting adversarial manipulation.

In addition, the project maintains interpretable diagnostics (per-client residual statistics and weighting statistics) to support auditing and to identify systematic patterns of misclassification that could disproportionately affect particular client populations. These practices contribute to a more transparent and accountable engineering process for systems intended to operate across diverse participants.

\section{Quality Management}

Quality was maintained using standard software and experimental practice. All code and configurations were tracked with version control to support traceability and reproducibility. Experiments were run with consistent settings (including fixed random seeds where appropriate) and key metrics were logged in a structured way (e.g., accuracy, $\tau/\mu/\nu$ statistics, aggregation weights, and attacker-aware diagnostics).

Basic validation checks were built into the pipeline, including numerical stability safeguards and sanity checks on intermediate values (such as weight non-negativity and normalisation). We also ran benign-condition regression tests to ensure the defence does not degrade normal training, and used baseline comparisons and ablations to verify the effect of each component.

\section{Safety Plan}

This project is a software-only study on public benchmark datasets and does not involve laboratory
equipment, hardware prototyping, human participants, or hazardous materials. Therefore, there are no
significant physical safety risks.

The main remaining risks are minimal and relate to good computing practice: ensuring code and results are
stored securely, avoiding any sensitive or personally identifiable data in datasets/logs, and documenting
the experimental setup to prevent misuse or misinterpretation of results. Any future deployment in
safety-critical domains would require a separate, domain-specific safety assessment.
